%% ========================================================================
%% Brāhmasphuṭasiddhānta Part 2 — Chunk 3 (Pages 1119–1258)
%% BILINGUAL English-Sanskrit Edition
%% LEFT: Sanskrit in Devanagari  |  RIGHT: English translation with IAST
%% ========================================================================

\documentclass[11pt,a4paper]{article}

%% -------- Core Packages --------
\usepackage{fontspec}
\usepackage{polyglossia}
\setmainlanguage{english}
\setotherlanguage{sanskrit}

\usepackage{amsmath,amssymb,amsthm}
\usepackage{geometry}
\usepackage{graphicx}
\usepackage{xcolor}
\usepackage{titlesec}
\usepackage{fancyhdr}
\usepackage{enumitem}
\usepackage{array,booktabs,longtable,multirow}
\usepackage{hyperref}
% \usepackage{verse}  % Package corrupted (404 HTML) -- not needed
\usepackage{lettrine}
\usepackage{microtype}
\usepackage{tocloft}
\usepackage{indentfirst}
\usepackage{paracol}
\usepackage{multicol}

%% -------- Page Geometry --------
\geometry{
  paperwidth=210mm,paperheight=297mm,
  left=18mm,right=18mm,top=25mm,bottom=25mm,
  headheight=14pt,footskip=30pt
}

%% -------- Font Configuration --------
\setmainfont[
  Path=/usr/share/fonts/opentype/linux-libertine/,
  BoldFont=LinLibertine_RB.otf,
  ItalicFont=LinLibertine_RI.otf,
  BoldItalicFont=LinLibertine_RBI.otf]{LinLibertine_R.otf}
\setsansfont[
  Path=/usr/share/fonts/opentype/linux-libertine/,
  BoldFont=LinBiolinum_RB.otf,
  ItalicFont=LinBiolinum_RI.otf]{LinBiolinum_R.otf}
\setmonofont[
  Path=/usr/share/fonts/truetype/dejavu/,
  Scale=MatchLowercase
]{DejaVuSansMono.ttf}

%% Devanagari (Sanskrit/Hindi) Font
\newfontfamily\devanagarifont[Script=Devanagari,ExternalLocation]{/usr/share/fonts/truetype/noto/NotoSerifDevanagari-Regular.ttf}
\newfontfamily\devanagarifontbf[Script=Devanagari,ExternalLocation]{/usr/share/fonts/truetype/noto/NotoSerifDevanagari-Bold.ttf}

%% -------- Colors --------
\definecolor{titlecolor}{RGB}{45,55,72}
\definecolor{sectioncolor}{RGB}{55,65,81}
\definecolor{notecolor}{RGB}{120,53,15}
\definecolor{rulecolor}{RGB}{180,160,140}
\definecolor{versebg}{RGB}{250,248,245}
\definecolor{sanskritbg}{RGB}{255,253,250}
\definecolor{englishbg}{RGB}{252,252,255}
\definecolor{columnsep}{RGB}{160,140,120}

%% -------- Section Formatting --------
\titleformat{\section}
  {\normalfont\Large\bfseries\color{sectioncolor}}
  {\thesection}{1em}{}
\titleformat{\subsection}
  {\normalfont\large\bfseries\color{sectioncolor}}
  {\thesubsection}{1em}{}
\titleformat{\subsubsection}
  {\normalfont\normalsize\bfseries\color{sectioncolor}}
  {\thesubsubsection}{1em}{}

%% -------- Header/Footer --------
\pagestyle{fancy}
\fancyhf{}
\fancyhead[L]{\small\textit{Brāhmasphuṭasiddhānta — Bilingual Edition — Part II, Chunk 3}}
\fancyhead[R]{\small\thepage}
\renewcommand{\headrulewidth}{0.4pt}
\renewcommand{\footrulewidth}{0pt}

%% -------- Custom Commands --------
\newcommand{\dev}[1]{{\devanagarifont #1}}
\newcommand{\devbf}[1]{{\devanagarifontbf #1}}

%% Parallel column commands
\newcommand{\sktcol}[1]{\fcolorbox{rulecolor}{sanskritbg}{\parbox{\linewidth}{\dev{#1}}}}
\newcommand{\engcol}[1]{\fcolorbox{rulecolor}{englishbg}{\parbox{\linewidth}{#1}}}

%% Verse environment for Sanskrit
\newenvironment{sanskritverse}{%
  \begin{center}\begin{minipage}{0.95\textwidth}
  \setlength{\leftskip}{1em}
  \setlength{\rightskip}{1em}
  \dev
}{\end{minipage}\end{center}}

%% Parallel verse environment
\newenvironment{paralleltext}{%
  \begin{paracol}{2}
  \backgroundcolor{c[0](.95\columnsep,.95\textheight)}{sanskritbg}
  \backgroundcolor{c[1](.95\columnsep,.95\textheight)}{englishbg}
  \def\sanskritcolumn{\switchcolumn[0]*}
  \def\englishcolumn{\switchcolumn[1]*}
}{\end{paracol}}

%% -------- Paracol Settings --------
\columnsep{10pt}
\setlength{\columnseprule}{0.4pt}
\def\columnseprulecolor{\color{columnsep}}

%% -------- Hyperref Setup --------
\hypersetup{
  colorlinks=true,
  linkcolor=sectioncolor,
  citecolor=sectioncolor,
  urlcolor=sectioncolor,
  pdfencoding=auto,
  pdftitle={Brāhmasphuṭasiddhānta Part 2 Chunk 3 — Bilingual Edition},
  pdfauthor={Brahmagupta, with commentary by Sudhākara Dvivedī}
}

%% ========================================================================
\newcommand{\vs}{\vspace{0.5em}}
\begin{document}

%% ========================================================================
%% TITLE PAGE
%% ========================================================================
\begin{titlepage}
\begin{center}
\vspace*{1.5cm}

{\devanagarifont\fontsize{26}{30}\selectfont\color{titlecolor} ब्राह्मस्फुटसिद्धान्तः}\\[0.8cm]
{\fontsize{20}{24}\selectfont\color{sectioncolor} Brāhmasphuṭasiddhānta}\\[1.2cm]

\rule{0.7\textwidth}{1pt}\\[0.8cm]

{\LARGE\textbf{BILINGUAL EDITION}}\\[0.3cm]
{\large Sanskrit (Devanagari) $\parallel$ English (with IAST)}\\[1.2cm]

{\Large\textbf{Part Two — Chunk 3}}\\[0.3cm]
{\large (Pages 281--420 of PDF $\approx$ Original pages 1119--1258)}\\[1.5cm]

{\large By \textbf{Brahmagupta} (6th century CE)}\\[0.5cm]
{\large With Commentary (\dev{नूतनतिलक}) by}\\[0.2cm]
{\large \textbf{Śrī Kṛpāludatta Sūnu Sudhākara Dvivedī}}\\[1.5cm]

\rule{0.7\textwidth}{0.4pt}\\[0.8cm]

{\normalsize Contains the following chapters:}\\[0.3cm]
\begin{minipage}{0.8\textwidth}
\begin{center}
{\devanagarifont चन्द्रग्रहणविकारः} — Lunar Eclipse Variations\\[0.15cm]
{\devanagarifont सूर्यग्रहणाधिकारः} — Solar Eclipse Chapter\\[0.15cm]
{\devanagarifont चन्द्रच्छायाधिकारः} — Moon's Shadow Chapter\\[0.15cm]
{\devanagarifont ग्रहोदयास्ताधिकारः} — Planetary Rising and Setting\\[0.15cm]
{\devanagarifont शुक्लोन्नत्यधिकारः} — Brightness of the Moon\\[0.15cm]
{\devanagarifont चन्द्रशुक्लोन्नत्यधिकारः} — Moon's Bright Elevation\\[0.15cm]
\end{center}
\end{minipage}

\vfill
{\small Typeset with XeLaTeX using Noto Serif Devanagari}\\
{\small Bilingual edition: Sanskrit left — English right}

\end{center}
\end{titlepage}

%% ========================================================================
%% TABLE OF CONTENTS
%% ========================================================================
\tableofcontents
\newpage

%% ========================================================================
%% EDITORIAL PREFACE TO THE BILINGUAL EDITION
%% ========================================================================
\section*{Editorial Preface}
\addcontentsline{toc}{section}{Editorial Preface}

This bilingual edition presents the Sanskrit text of Brāhmasphuṭasiddhānta Part II, Chunk 3, with facing English translations. The layout follows the standard format for critical editions of Sanskrit scientific texts:

\begin{itemize}
  \item \textbf{Left column:} Sanskrit text in Devanagari script, as found in the source edition.
  \item \textbf{Right column:} English translation with IAST romanization of technical terms.
\end{itemize}

The commentary (\dev{टीका}) by Sudhākara Dvivedī follows the traditional Indian \emph{pañcāṅga} structure of:
\begin{enumerate}
  \item \textbf{Sūtra} (\dev{सूत्रम्}) — The original verse in āryā or anuṣṭubh metre.
  \item \textbf{Sūtrārtha} (\dev{सूत्रार्थः}) — Literal meaning of the verse.
  \item \textbf{Hindi Bhāṣya} (\dev{हिन्दी भाष्य}) — Explanation in Hindi.
  \item \textbf{Upapatti} (\dev{उपपत्तिः}) — Mathematical proof or demonstration.
\end{enumerate}

\medskip
\noindent\rule{\textwidth}{0.4pt}

\newpage


%% ========================================================================
%% CHAPTER 1: CHANDRAGRAHAṆAVIKĀRAḤ (Lunar Eclipse Variations)
%% ========================================================================
\section{\texorpdfstring{{\devanagarifont चन्द्रग्रहणविकारः} — Candragrahaṇavikāraḥ (Lunar Eclipse Variations)}{Candragrahaṇavikāraḥ (Lunar Eclipse Variations)}}
\label{sec:chandra-grahana}

\begin{center}
\textcolor{rulecolor}{\rule{0.8\textwidth}{0.5pt}}
\end{center}

\subsection*{Introduction}

This chapter (\dev{चन्द्रग्रहणविकारः}) deals with the computation of lunar eclipses, following Brahmagupta's exposition in the \dev{ब्राह्मस्फुटसिद्धान्त}. The commentary (\dev{सुधाकर-टीका}) by Paṇḍita Kṛpāludatta Sūnu Sudhākara Dvivedī provides elaborate explanations in Hindi, with mathematical demonstrations (\dev{उपपत्ति}).

\begin{center}
\textcolor{rulecolor}{\rule{0.8\textwidth}{0.5pt}}
\end{center}

%% ========================================================================
\subsection{Verses 11--12: Computation of \texorpdfstring{{\devanagarifont इष्टग्रास}}{iṣṭagrāsa} (Desired Eclipse Magnitude)}

\vspace{0.5em}
\noindent\textbf{Sūtra (\dev{सूत्रम्}) — Verses 11--12:}

\begin{paracol}{2}
\backgroundcolor{c[0](.95\columnsep,.95\textheight)}{sanskritbg}
\backgroundcolor{c[1](.95\columnsep,.95\textheight)}{englishbg}

\begin{sanskritverse}
इष्टन्यूनस्थितिदलगुणा भुक्तिविल्लेपलिप्ता पश्चादष्टा भवति हि भुजः कोटिरिष्टेणुवाणः \\
तद्वेकोद्भवमपि पदं कर्णं एतेन हीनं मानैक्यार्धं स्फुटमिह भवेदाकलित छन्नमानम्॥ ११--१२ ॥
\end{sanskritverse}

\switchcolumn

\textbf{Translation:} ``The half-duration (\textit{sthitidala}) diminished by the desired quantity, multiplied by the daily motion (\textit{bhukti}) and divided by the deflection (\textit{villepa}), gives the \textit{bhuja} (sine). The \textit{koṭi} (cosine) is computed by the rule of the desired quantity (\textit{iṣṭeṇu}). From the half-sum of the diameters (\textit{mānaikyārdha}), the \textit{karṇa} (hypotenuse) being subtracted, the result gives the true obscuration (\textit{channa-māna}).''

\end{paracol}

\vspace{0.5em}

\begin{paracol}{2}
\backgroundcolor{c[0](.95\columnsep,.95\textheight)}{sanskritbg}
\backgroundcolor{c[1](.95\columnsep,.95\textheight)}{englishbg}

\noindent\textbf{Sūtrārtha (\dev{सूत्रार्थः}):}
\begin{sanskritverse}
इष्टन्यूनस्थितिदलगुणा भुक्तिविल्लेपलिप्ता पश्चादष्टा भवति हि भुजः कोटिरिष्टेणुवाणः।
तद्वेकोद्भवमपि पदं कर्णं एतेन हीनं मानैक्यार्धं स्फुटमिह भवेदाकलित छन्नमानम्॥
\end{sanskritverse}

\switchcolumn

\textbf{Literal Meaning:} The verse gives the method for computing the \textit{iṣṭagrāsa} (desired eclipse magnitude). From the \textit{sthitidala} (half-duration) diminished by the desired quantity, multiplied by the \textit{bhukti} (daily motion) and divided by the \textit{villepa} (deflection), one gets the \textit{bhuja} (sine). The \textit{koṭi} (cosine) is computed by the rule of \textit{iṣṭeṇu} (desired quantity). From the \textit{mānaikyārdha} (half-sum of diameters), the \textit{karṇa} (hypotenuse) being subtracted, the result gives the true obscuration.

\end{paracol}

\vspace{0.5em}

\begin{paracol}{2}
\backgroundcolor{c[0](.95\columnsep,.95\textheight)}{sanskritbg}
\backgroundcolor{c[1](.95\columnsep,.95\textheight)}{englishbg}

\noindent\textbf{Hindi Bhāṣya (\dev{हिन्दी भाष्य}):}
\begin{sanskritverse}
इष्ट रहित स्थितिदल को भुक्ति से गुणा कर विल्लेप से भाग देने से भुज होता है, इस फल करके
सूर्य, चन्द्र और पात को चालन देकर तात्कालिक चन्द्रशर साधन करना वह कोटि है। उन दोनों
(भुज और कोटि) का वर्गयोग मूल कर्ण होता है, मानैक्यार्ध (छाद्य और छादक बिम्ब के व्यासार्ध
योग) में से कर्ण को घटाने से जो शेष रहता है वह इष्टग्रास होता है इति॥ ११--१२ ॥
\end{sanskritverse}

\switchcolumn

\textbf{English Translation of Hindi Commentary:} ``Subtracting the desired (\textit{iṣṭa}) from the \textit{sthitidala}, multiplying by the \textit{bhukti} and dividing by the \textit{villepa}, one obtains the \textit{bhuja}. By applying this result to the motions of the Sun, Moon, and node (\textit{pāta}), one computes the instantaneous lunar latitude (\textit{tātkālika candra-śara}), which is the \textit{koṭi}. The square root of the sum of their squares is the \textit{karṇa}. From the \textit{mānaikyārdha} (sum of the radii of the eclipsed and eclipsing bodies), subtracting the \textit{karṇa}, the remainder is the desired magnitude (\textit{iṣṭagrāsa}).''

\end{paracol}

\vspace{0.5em}

\begin{paracol}{2}
\backgroundcolor{c[0](.95\columnsep,.95\textheight)}{sanskritbg}
\backgroundcolor{c[1](.95\columnsep,.95\textheight)}{englishbg}

\noindent\textbf{Upapatti (\dev{उपपत्तिः}) — Proof:}
\begin{sanskritverse}
स्पर्शी के बाद जितनी इष्ट घटी में ग्रास-ज्ञान अपेक्षित हो उस इष्ट घटी को स्थित्यर्ध में से घटा
देना, शेष से अनुपात करते हैं। यदि साठ घटी में रवि और चन्द्र की गत्यन्तर कला पाते हैं तो
इष्टोन स्थित्यर्ध घटी में क्या इससे जो फल होता है वह भुज है, अनुपातागत फल से चालित रवि,
चन्द्र और पात से तात्कालिक चन्द्रवार साधन करना वह कोटि है। इन दोनों (भुज और कोटि)
का वर्गयोग मूल कर्ण होता है, मानैक्यार्ध में से कर्ण को घटाने से इष्ट ग्रास होता है।
\end{sanskritverse}

\switchcolumn

\textbf{English Translation of Proof:} ``After the contact (\textit{sparśī}), for whatever desired number of ghaṭīs the knowledge of the \textit{grāsa} is required, that desired number of ghaṭīs should be subtracted from the \textit{sthityardha}. The remainder is used in proportion. If in sixty ghaṭīs the difference in motion between the Sun and Moon gives a certain number of minutes (\textit{kalā}), then [by proportion] in the \textit{iṣṭa}-diminished \textit{sthityardha} ghaṭīs, whatever result is obtained is the \textit{bhuja}. From the result obtained by proportion, by applying the motions of the Sun, Moon, and node, the instantaneous lunar latitude (\textit{tātkālika candra-vāra}) is computed --- that is the \textit{koṭi}. The square root of the sum of the squares of these two (\textit{bhuja} and \textit{koṭi}) is the \textit{karṇa}. Subtracting the \textit{karṇa} from the \textit{mānaikyārdha} gives the desired \textit{grāsa}.''

\end{paracol}

%% ========================================================================
\subsection{Verses 13--14: Computation of \texorpdfstring{{\devanagarifont इष्टग्रासात्काल}}{iṣṭagrāsāt-kāla} (Time from Desired Magnitude)}

\vspace{0.5em}
\noindent\textbf{Sūtra (\dev{सूत्रम्}) — Verses 13--14:}

\begin{paracol}{2}
\backgroundcolor{c[0](.95\columnsep,.95\textheight)}{sanskritbg}
\backgroundcolor{c[1](.95\columnsep,.95\textheight)}{englishbg}

\begin{sanskritverse}
प्रसकृत्-ग्रासकलोनप्रमाणयुतिदलकृते विशोध्य कृतिम्।
तात्कालिकविक्षेपस्य शेषमूलं कृतं तिथिवत्॥ १३ ॥
\end{sanskritverse}

\begin{sanskritverse}
प्रग्रहणस्थित्यर्धात् प्रोह्य प्रग्रहणतो भवेत् कालः।
सौक्षं विशोष्य मोक्षस्थित्यर्धात् प्राग् भवेन्मोक्षात्॥ १४ ॥
\end{sanskritverse}

\switchcolumn

\textbf{Translation — Verse 13:} ``From the square of the half-sum of the measure diminished by the \textit{grāsakalā}, having subtracted the square of the instantaneous latitude (\textit{tātkālika-vikṣepa}), the square root of the remainder --- made like a \textit{tithi} [i.e., multiplied by 60 and divided by the difference of motions] --- [gives a time].''

\textbf{Translation — Verse 14:} ``Having subtracted [this time] from the \textit{pragrahaṇa-sthit yardha}, the time from the \textit{pragrahaṇa} [i.e., from contact] is obtained. Having subtracted from the \textit{mokṣa-sthit yardha}, [one obtains the time] before the \textit{mokṣa} [release].''

\end{paracol}

\vspace{0.5em}

\begin{paracol}{2}
\backgroundcolor{c[0](.95\columnsep,.95\textheight)}{sanskritbg}
\backgroundcolor{c[1](.95\columnsep,.95\textheight)}{englishbg}

\noindent\textbf{Sūtrārtha (\dev{सूत्रार्थः}):}
\begin{sanskritverse}
ग्रासकलोनप्रमाणयुतिदलकृते ग्रासोनमानैक्यार्धवर्गात् तात्कालिकशरस्य कृति विद्धोध्य
शेषशूलं तिथिवदसकृत् कृतम्। अर्थात् शेषमूलं षष्टिगुणां रविचन्द्रगत्यन्तरहतं फलकालेन
रविचन्द्रपातानु प्रचाल्य तात्कालिकविक्षेपं संसाध्य तस्मात् पुनः ग्रासिकलोनप्रमाणेत्यादिना
असकृत् शेषमूलकालः स्थिरीभवति तं प्रग्रहणस्थित्यर्धात् (स्पर्शिकस्थित्यर्धात्) प्रोह्य
(हित्वा) प्रग्रहणतः (स्पर्शादनन्तरं) कालो भवेत्।
\end{sanskritverse}

\switchcolumn

\textbf{Literal Meaning:} ``From the square of the half-sum of the measure [of diameters] diminished by the \textit{grāsakalā}, having subtracted the square of the instantaneous \textit{śara} (latitude), the square root of the remainder --- made \textit{tithi}-like (\textit{tithivat}) [by multiplying by 60 and dividing by the difference of motions of Sun and Moon] --- is the result. That is, having multiplied the square root of the remainder by 60, divided by the difference of motions of Sun and Moon, and having applied the motions of Sun, Moon, and node by that resulting time, the instantaneous latitude (\textit{tātkālika-vikṣepa}) is computed. From that, again by the rule `\textit{grāsakalona-pramāṇa}', etc., iteratively (\textit{asakṛt}), the square-root-time becomes stable. That [time], subtracted from the \textit{pragrahaṇa-sthit yardha} (contact half-duration), the time from the \textit{pragrahaṇa} (from contact) is obtained.''

\end{paracol}

\vspace{0.5em}

\begin{paracol}{2}
\backgroundcolor{c[0](.95\columnsep,.95\textheight)}{sanskritbg}
\backgroundcolor{c[1](.95\columnsep,.95\textheight)}{englishbg}

\noindent\textbf{Hindi Bhāṣya (\dev{हिन्दी भाष्य}):}
\begin{sanskritverse}
ग्रासकलोनप्रमाणयुतिदलकृतेः (ग्रासरहितमानैक्यार्धवर्गात्) तात्कालिकशरस्य कृति
(वर्ग) विशोध्य (हित्वा) शेषस्थ मूल तिथिवदसकृत् कृतमर्थात् शेषमूलं षष्टिगुणं
रविचन्द्रयोगंत्यन्तरभक्तं फलकालेन रविचन्द्रपातानु प्रचाल्य तात्कालिकचन्द्रशरं संसाध्य
तस्मात् पुनर्ग्रासकलोनप्रमाणयुतिदलकृतेरित्यादिना असकृत् शेषमूलं कालः स्थिरीभवति
ते प्रग्रहणस्थित्यर्धात् (स्पर्शिकस्थित्यर्धात्) प्रोह्य (हित्वा) प्रग्रहणतः (स्पर्शादनन्तरं)
कालो भवेदर्थात् स्पर्शानन्तरमेतावतीष्ठकाले तावानिष्टग्रासो भवति॥ १३--१४ ॥
\end{sanskritverse}

\switchcolumn

\textbf{English Translation:} ``From the square of the \textit{grāsa}-diminished \textit{mānaikyārdha}, having subtracted the square of the instantaneous \textit{śara}, the square root of the remainder --- made \textit{tithi}-like iteratively --- that is, having multiplied the square root of the remainder by 60, divided by the combined difference of motions of Sun and Moon, and having applied the motions of Sun, Moon, and node by that resulting time, the instantaneous lunar latitude is computed. From that, again by the rule beginning `\textit{grāsakalona-pramāṇa...}', iteratively (\textit{asakṛt}), the square-root-time becomes stable. That, subtracted from the contact half-duration (\textit{pragrahaṇa-sthit yardha}), gives the time from contact, meaning that after contact, in that desired time, the desired \textit{grāsa} is obtained.''

\end{paracol}

\vspace{0.5em}

\begin{paracol}{2}
\backgroundcolor{c[0](.95\columnsep,.95\textheight)}{sanskritbg}
\backgroundcolor{c[1](.95\columnsep,.95\textheight)}{englishbg}

\noindent\textbf{Upapatti (\dev{उपपत्तिः}):}
\begin{sanskritverse}
तत्कालशरस्याज्ञानान्मध्यकालिकशरात् कर्म कृतमतोऽसकृद्धिघिना स्फुटकालसाधनमुचितम्।
शेषोपपत्ति ``ग्रसोन मानैक्यदलस्य वर्गाद् विक्षेपकृत्या रहिता'' इत्यादिभास्करविधिना
स्फुटा ॥ १३--१४ ॥
\end{sanskritverse}

\switchcolumn

\textbf{English Translation of Proof:} ``Due to the ignorance of the latitude at that time (\textit{tātkāla}), the computation has been done using the middle-time latitude. Therefore, it is proper to determine the accurate time by an iterative method (\textit{asakṛt-ghātinā}). The remaining proof [follows] the method of Bhāskara: `From the square of half the sum of diameters diminished by the \textit{grāsa}, having subtracted the square of the \textit{vikṣepa}...' etc.''

\end{paracol}

%% ========================================================================
\subsection{Verse 9: Sphuṭīkaraṇa of \texorpdfstring{{\devanagarifont स्थितिविमर्दार्ध}}{sthitivimardārdha} (Refinement of Half-Duration of Totality)}

\vspace{0.5em}
\noindent\textbf{Sūtra (\dev{सूत्रम्}) — Verse 9:}

\begin{paracol}{2}
\backgroundcolor{c[0](.95\columnsep,.95\textheight)}{sanskritbg}
\backgroundcolor{c[1](.95\columnsep,.95\textheight)}{englishbg}

\begin{sanskritverse}
षष्ट्या विभाजिता स्थितिविमर्दार्धनाडिकागुणा स्वगतिः।
छाद्ये रवीन्दुपातेष्वुभयसकृत् तैर्धनमन्ते॥ ९ ॥
\end{sanskritverse}

\switchcolumn

\textbf{Translation:} ``The own motion (\textit{svagati}) --- [i.e.] the motions of the Sun, Moon, and node --- multiplied by the \textit{sthitivimardārdha} nāḍikās and divided by sixty; both [motions] are applied [to the Sun, Moon, and node] repeatedly in the eclipsed body (\textit{chādya}), [and] the result is additive or subtractive.''

\end{paracol}

\vspace{0.5em}

\begin{paracol}{2}
\backgroundcolor{c[0](.95\columnsep,.95\textheight)}{sanskritbg}
\backgroundcolor{c[1](.95\columnsep,.95\textheight)}{englishbg}

\noindent\textbf{Sūtrārtha (\dev{सूत्रार्थः}):}
\begin{sanskritverse}
स्वगतिः (रविगतिः, चन्द्रगतिः, पातगतिश्च) स्थितिविमर्दार्धनाडिका गुणा
(स्थित्यर्ध घटीमिविमर्दार्ध घटीभिः पृथक् पृथक् गुणिता) षष्ट्या विभाजिता छाद्ये
रविचन्द्रपातेषु उभयसकृत् तैर्धनमन्ते।
\end{sanskritverse}

\switchcolumn

\textbf{Literal Meaning:} ``The own motion --- the motion of the Sun, the motion of the Moon, and the motion of the node --- multiplied by the \textit{sthitivimardārdha} nāḍikās (the half-duration of totality ghaṭīs multiplied separately) and divided by sixty; applied to the Sun, Moon, and node in the eclipsed body, both [additive and subtractive] repeatedly, the result being additive or subtractive.''

\end{paracol}

\vspace{0.5em}

\begin{paracol}{2}
\backgroundcolor{c[0](.95\columnsep,.95\textheight)}{sanskritbg}
\backgroundcolor{c[1](.95\columnsep,.95\textheight)}{englishbg}

\noindent\textbf{Upapatti (\dev{उपपत्तिः}):}
\begin{sanskritverse}
``स्थित्यर्धनाडीगुणिता स्वमुक्ति'' इत्यादिना एवं विमर्दार्धफलोनयुक्त इत्यादिना च
भास्करविधिना स्फुटा ॥ ९ ॥
\end{sanskritverse}

\switchcolumn

\textbf{English Translation:} ``The [method is] made accurate by the method of Bhāskara: `\textit{sthit yardha-nāḍī-guṇitā sva-mukti}' etc., and `\textit{vimardārdha-phala-ona-yukta}' etc.''

\end{paracol}

\newpage


%% ========================================================================
%% CHAPTER 2: SŪRYAGRAHAṆĀDHIKĀRAḤ (Solar Eclipse Chapter)
%% ========================================================================
\section{\texorpdfstring{{\devanagarifont सूर्यग्रहणाधिकारः} — Sūryagrahaṇādhikāraḥ (Solar Eclipse Chapter)}{Sūryagrahaṇādhikāraḥ (Solar Eclipse Chapter)}}
\label{sec:surya-grahana}

\begin{center}
\textcolor{rulecolor}{\rule{0.8\textwidth}{0.5pt}}
\end{center}

This chapter (\dev{सूर्यग्रहणाधिकारः}, the fifth chapter of the text) deals with the computation of solar eclipses. The chapter begins with the statement of purpose (\dev{प्रयोजन}) and proceeds through the computation of the five \dev{ज्या} (sines) required for eclipse calculation: \dev{उदयज्या}, \dev{मध्यज्या}, \dev{रविदयाकर्ण}, \dev{विशिमशङ्कु}, and \dev{हृक्षेप}.

%% ========================================================================
\subsection{Verse 1: Statement of Purpose}

\vspace{0.5em}
\noindent\textbf{Sūtra (\dev{सूत्रम्}) — Verse 1:}

\begin{paracol}{2}
\backgroundcolor{c[0](.95\columnsep,.95\textheight)}{sanskritbg}
\backgroundcolor{c[1](.95\columnsep,.95\textheight)}{englishbg}

\begin{sanskritverse}
अथ सूर्यग्रहणाधिकारः प्रारभ्यते। तत्रादौ तदारम्भप्रयोजनमाह—
\end{sanskritverse}

\begin{sanskritverse}
ग्रहभास्करान्तरं प्राक् पद बाधकं ग्रहान्तरे यस्मात्।
स्वांशैर्यथाह्यास्तस्माद् द्युते तदानयनम्॥ १ ॥
\end{sanskritverse}

\switchcolumn

\textbf{Introduction:} ``Now the chapter on solar eclipse is begun. First, the purpose of its commencement is stated ---''

\textbf{Translation:} ``Because the inter-planetary [distance] computed by the \textit{grahabhāskara} [method] previously is an obstruction in the inter-planetary [computation], therefore I shall state the computation of \textit{dyuti} (parallax) with its own degrees, so that there may be agreement.''

\textit{Technical note:} The \textit{grahabhāskara} refers to the method given by earlier astronomers (like Bhāskara I) for computing the distance between planets, which Brahmagupta finds to be defective for solar eclipse computation.

\end{paracol}

\vspace{0.5em}

\begin{paracol}{2}
\backgroundcolor{c[0](.95\columnsep,.95\textheight)}{sanskritbg}
\backgroundcolor{c[1](.95\columnsep,.95\textheight)}{englishbg}

\noindent\textbf{Sūtrārtha (\dev{सूत्रार्थः}):}
\begin{sanskritverse}
प्राक् पूर्वे क्षितिजे प्रभास्करान्तरैः स्वांशैः स्वकालांशैः पश्चात् पश्चिमक्षितिजे
ग्रहास्तरैः स्वकालांशैर्यथासंख्यं ग्रहाः हश्याहश्याः भवन्ति। तस्मात् तदानयनं
प्रहर्कान्तरानयनं वक्ष्ये इति। अर्थात् प्राक्-क्षितिजे यस्मिन् दिने ग्रहोदयादनन्तरं
कालांशघटीभिः रवेरुदयस्तस्मिन् दिने ग्रहो रात्रिशेषे हश्यो भवति। एवं यस्मिन् दिने
पश्चिमक्षितिजे रव्यस्तानन्तरं कालांशघटीमित्रं हस्यास्तमयस्तस्मिन् दिने सायंकाले
ग्रहस्य हश्यत्वम्॥ १ ॥
\end{sanskritverse}

\switchcolumn

\textbf{Literal Meaning:} ``Previously, in the eastern horizon (\textit{pūrve kṣitije}), by the degrees of the \textit{grahabhāskara} method, the planets become visible or invisible (\textit{haśya-āhaśya}) in their own time-degrees. Later, in the western horizon (\textit{paścima-kṣitije}), by the degrees of the planetary distances, [the same] according to number. Therefore, I shall state that computation --- the computation of the \textit{praharāntara}. That is: in the eastern horizon, on whatever day after the rising of the planet, the rising of the Sun [occurs] after a certain number of time-ghaṭīs, on that day the planet becomes invisible (\textit{haśya}) in the remaining night. Similarly, in the western horizon, after the setting of the Sun, when the setting [occurs] after a certain number of time-ghaṭīs, on that day in the evening the planet's invisibility (\textit{haśyatva}) [is determined].''

\end{paracol}

\vspace{0.5em}

\begin{paracol}{2}
\backgroundcolor{c[0](.95\columnsep,.95\textheight)}{sanskritbg}
\backgroundcolor{c[1](.95\columnsep,.95\textheight)}{englishbg}

\noindent\textbf{Hindi Bhāṣya (\dev{हिन्दी भाष्य}):}
\begin{sanskritverse}
जिस कारण से सूर्यग्रहण में पञ्चज्या (उदयज्या, मध्यज्या, रविदयाकर्ण, विशिमशङ्कु,
हृक्षेप) से भास्करादि आचार्यों ने जो प्रानयन किया है उससे गणितैक्य (वेधोपलब्ध और
गणितागत विषयों की तुल्यता) नहीं होता है उस कारण से जैसे उनका ऐक्य (वेधोपलब्ध और
गणितागत विषयों के समत्व) हो वैसे मैं प्रानयन को कहता हूँ। तिथ्यन्ते इसका ग्रहे से
सम्बन्ध है, उदयज्या -- भ्रमा, मध्यज्या -- दशमलग्ननतांशज्या, विशिमशङ्कु -- हग्गति,
हृक्षेप -- विशिमनतांशज्या, सूर्यसिद्धान्त में भी पूर्वोक्त पञ्चज्या ही से सूर्यग्रहण का
प्रानयन है, पौलिशसिद्धान्त में सूर्यग्रहण में चन्द्र की पृथक् पञ्चज्या साधित है इसलिए
वहाँ द्विज्या विधान से सूर्यग्रहण उपनिबद्ध है इति ॥ १ ॥
\end{sanskritverse}

\switchcolumn

\textbf{English Translation:} ``For whatever reason, in solar eclipses, the five \textit{jyā}s (\textit{udayajyā}, \textit{madhyajyā}, \textit{ravidayākarṇa}, \textit{viśimaśaṅku}, \textit{hṛkṣepa}) --- by which earlier teachers like Bhāskara computed the \textit{prāṇayana} --- from that, the \textit{gaṇitaikya} (agreement between observed and computed values) is not obtained. For that reason, so that their agreement may exist, I state the \textit{prāṇayana} [method]. The \textit{tithi-anta} [refers to] the relation with the \textit{graha}: \textit{udayajyā} [refers to] parallax (\textit{bhramā}), \textit{madhyajyā} [to] the sine of the zenith-distance of the tenth degree (\textit{daśama-lagna-natāṃśajyā}), \textit{viśimaśaṅku} [to] the ascensional difference (\textit{hagati}), \textit{hṛkṣepa} [to] the sine of the zenith-distance of the nonagesimal (\textit{viśima-natāṃśajyā}). In the \textit{Sūryasiddhānta} too, the solar eclipse \textit{prāṇayana} is by the aforesaid five \textit{jyā}s. In the \textit{Pauliśasiddhānta}, in solar eclipses, separate five \textit{jyā}s of the Moon are computed; therefore there the solar eclipse is accomplished by the two-\textit{jyā} method.''

\end{paracol}

%% ========================================================================
\subsection{Verses 2--3: Conditions for \texorpdfstring{{\devanagarifont लम्बन}}{lambana} (Deflection) and \texorpdfstring{{\devanagarifont अवनति}}{avanati} (Depression)}

\vspace{0.5em}
\noindent\textbf{Sūtra (\dev{सूत्रम्}) — Verses 2--3:}

\begin{paracol}{2}
\backgroundcolor{c[0](.95\columnsep,.95\textheight)}{sanskritbg}
\backgroundcolor{c[1](.95\columnsep,.95\textheight)}{englishbg}

\begin{sanskritverse}
विश्रिमलग्नसमं एकं न लम्बनं तदधिकोनके भवति।
तस्य क्रान्तिज्योदक् यदा अक्षनोवा समा न तदा॥ २ ॥
\end{sanskritverse}

\begin{sanskritverse}
अवनतिरतोऽन्यथा भवति सम्भवे तदुदयैव लग्नसमम्।
तदुदितघटिकास्तच्छङ्कुस्तच्चरप्राणैः परे॥ ३ ॥
\end{sanskritverse}

\switchcolumn

\textbf{Translation — Verse 2:} ``When one [planet] is equal to the \textit{viśrima-lagna} (nonagesimal ascendant), there is no \textit{lambana} (deflection in longitude). When it is greater or less than that, \textit{lambana} occurs. When its declination-sine (\textit{krāntijyā}) is equal to the latitude-sine (\textit{akṣajyā}), then there is no \textit{avanati} (depression).''

\textbf{Translation — Verse 3:} ``Otherwise, \textit{avanati} occurs. When both \textit{lambana} and \textit{avanati} are possible, then the rising [point] equal to the \textit{lagna}, the \textit{ghaṭikās} elapsed since its rising, its \textit{śaṅku} (gnomon), and the \textit{prāṇa}s of its motion --- [are computed] thereafter.''

\end{paracol}

\vspace{0.5em}

\begin{paracol}{2}
\backgroundcolor{c[0](.95\columnsep,.95\textheight)}{sanskritbg}
\backgroundcolor{c[1](.95\columnsep,.95\textheight)}{englishbg}

\noindent\textbf{Sūtrārtha (\dev{सूत्रार्थः}):}
\begin{sanskritverse}
तिथ्यन्ते गणितागतदर्शान्तिकाले पूर्वविधिना त्रिप्रश्नोक्तेन लग्नं कृत्वा विश्रिमं
कार्यम्। तेन विश्रिमलग्नेन समेकं रवौ सति लम्बनं न भवति! तस्माद् विश्रिमादधिके
कोने रवौ लम्बनं भवतीत्यर्थत् एव सिध्यति। एवं तस्य विश्रिमस्योत्तरा क्रान्तिज्या यदा
स्वदेशाक्षजोवा समा तदा अवनतिर्न भवति अतोऽन्यथा चावनतिर्भवति।
लम्बनावनत्योः सम्भवे च तदुदयं स्वदेशीयराश्युदयैषि-लग्नसमं कृत्वा
आद्यूनस्य भोग्योद्धिकभुक्तियुक्तो मध्योदयाद्य इति त्रिप्रश्नोक्तविधिं कृत्वा
तदुदितघटिकास्तस्य विश्रिमस्य गता घटिकाः साध्याः॥ २--३ ॥
\end{sanskritverse}

\switchcolumn

\textbf{Literal Meaning:} ``At the end of the \textit{tithi}, at the time of computed conjunction (\textit{gaṇitāgata-darśāntikāle}), having computed the \textit{lagna} by the previous method stated in the \textit{Tripraśna}, the \textit{viśrima} [nonagesimal] is to be computed. When the Sun is equal to one [point] with that \textit{viśrimalagna}, there is no \textit{lambana}. When the Sun is greater or less than that \textit{viśrima}, \textit{lambana} occurs --- this meaning is self-evident. Similarly, when the northern declination-sine (\textit{krāntijyā}) of that \textit{viśrima} is equal to the latitude-sine (\textit{akṣajyā}) of one's own place, then there is no \textit{avanati}; otherwise \textit{avanati} occurs. When both \textit{lambana} and \textit{avanati} are possible, then having made the rising equal to the \textit{lagna} --- the \textit{rāśyudaya-eṣi-lagna} of one's own place --- by the method stated in the \textit{Tripraśna}: `the first diminished by the \textit{bhoga}, increased by the \textit{bhukti}, [gives] the middle \textit{udaya}', etc., the \textit{ghaṭikās} elapsed since the rising of that \textit{viśrima} are to be computed.''

\end{paracol}

%% ========================================================================
\subsection{Verse 7: {\devanagarifont स्पष्टदर्शान्तकालिक} Computation (Computation of True Conjunction Time)}

\vspace{0.5em}
\noindent\textbf{Sūtra ({\devanagarifont सूत्रम्}) — Verse 7:}

\begin{paracol}{2}
\backgroundcolor{c[0](.95\columnsep,.95\textheight)}{sanskritbg}
\backgroundcolor{c[1](.95\columnsep,.95\textheight)}{englishbg}

\begin{sanskritverse}
घटिकामिरगुणिताः पञ्चभिर्भक्ता लब्धकलामियंति लम्बनम्।
तदा गणितागतदर्शान्तकालिका रविचन्द्रपाता ऊनाः कार्याः॥ ७ ॥
\end{sanskritverse}

\switchcolumn

\textbf{Translation:} ``The \textit{ghaṭikā-miṣa} [of \textit{lambana}] multiplied by five and divided [by some number] yields the \textit{labdha-kalā} called \textit{lambana}. Then the Sun, Moon, and node at the computed conjunction time (\textit{gaṇitāgata-darśānta-kālika}) are to be diminished [or increased].''

\end{paracol}

\vspace{0.5em}

\begin{paracol}{2}
\backgroundcolor{c[0](.95\columnsep,.95\textheight)}{sanskritbg}
\backgroundcolor{c[1](.95\columnsep,.95\textheight)}{englishbg}

\noindent\textbf{Sūtrārtha (\dev{सूत्रार्थः}):}
\begin{sanskritverse}
घटिकामिरगुणिताः पञ्चभ्यः पञ्चज्या भक्ता लब्धकलामियंति लम्बनं तदा
गणितागतदर्शान्तकालिका रविचन्द्रपाता ऊनाः (हीनाः) कार्याः, यदि लम्बनं धनं तदा
लब्धकलाभिर्गणितदर्शान्तकालिका रविचन्द्रपाता युक्ताः कार्यास्तदा स्पष्टदर्शान्तकालिका
रविचन्द्रपाता भवन्तीति ॥ ७ ॥
\end{sanskritverse}

\switchcolumn

\textbf{Literal Meaning:} ``The \textit{ghaṭikā-miṣa} multiplied by five, divided by five [?\ \textit{pañcajyā}], the result is the \textit{labdha-kalā} called \textit{lambana}. Then the Sun, Moon, and node at the computed conjunction time are to be diminished (\textit{ūnāḥ}). If the \textit{lambana} is positive (\textit{dhana}), then the Sun, Moon, and node at the computed conjunction time are to be increased by the \textit{labdha-kalā}; then those become the true Sun, Moon, and node at the true conjunction time (\textit{spaṣṭa-darśānta-kālika}).''

\end{paracol}

\vspace{0.5em}

\begin{paracol}{2}
\backgroundcolor{c[0](.95\columnsep,.95\textheight)}{sanskritbg}
\backgroundcolor{c[1](.95\columnsep,.95\textheight)}{englishbg}

\noindent\textbf{Upapatti (\dev{उपपत्तिः}):}
\begin{sanskritverse}
गणितागतदर्शान्तकालिकरविचन्द्रपातेभ्यः शरादिकं सर्वमानेतव्यम्। ऋणात्मके
लम्बने ``पञ्चषष्टिकामिरः पृथक् पृथक् रविचन्द्रपातगतिकला लम्यन्ते तदा
लम्बनघटिकाभिः किमित्यनुपातेन'' लब्धकलाभिर्गणितागतदर्शान्तकालिका
रविचन्द्रपाता हीनााः कार्याः, धनात्मके लम्बने लब्धकलाभिस्ते युक्तास्तदा ते
स्पष्टदर्शान्तकालिका भवेयुरेवेति, सिद्धान्तशेखरे ``गतिहतां लम्बननाडिकामिर्विभज्य
पञ्चया फललिप्तिकामिः। रवीन्दुपाताः सहिता धनाख्ये विवर्जिताश्च क्षयलम्बने ते''
इत्यनेन श्रीपतिनाचार्योक्तानुरूपमेवोक्तमिति ॥ ७ ॥
\end{sanskritverse}

\switchcolumn

\textbf{English Translation of Proof:} ``From the Sun, Moon, and node at the computed conjunction time, all [quantities] beginning with the \textit{śara} [latitude] are to be computed. When the \textit{lambana} is negative (\textit{ṛṇātmaka}): `The minutes of motion of the Sun, Moon, and node are separately multiplied by the \textit{lambana}-ghaṭikā-miṣa} of sixty-five --- by proportion: `then by the \textit{lambana}-ghaṭikā, what [result]?}' --- by the \textit{labdha-kalā}s, the Sun, Moon, and node at the computed conjunction time are to be diminished. When the \textit{lambana} is positive (\textit{dhanātmaka}), they are increased by the \textit{labdha-kalā}s; then they become the true [quantities] at the true conjunction time. In the \textit{Siddhāntaśekhara}: `Having divided the motion-multiplied \textit{lambana}-nāḍī-miṣa by five, the resulting minutes... The Sun, Moon, and node are added in positive [\textit{lambana}] and subtracted in negative \textit{lambana}' --- by this, what was stated by Śrīpati-nācārya is here spoken in accordance.''

\end{paracol}

\newpage


%% ========================================================================
%% CHAPTER 3: CHANDRACCHĀYĀDHIKĀRAḤ (Moon's Shadow Chapter)
%% ========================================================================
\section{\texorpdfstring{{\devanagarifont चन्द्रच्छायाधिकारः} — Candracchāyādhikāraḥ (Moon's Shadow Chapter)}{Candracchāyādhikāraḥ (Moon's Shadow Chapter)}}
\label{sec:chandra-chhaya}

\begin{center}
\textcolor{rulecolor}{\rule{0.8\textwidth}{0.5pt}}
\end{center}

This chapter (\dev{चन्द्रच्छायाधिकारः}) deals with the computation of the Moon's shadow, which is essential for understanding lunar phenomena. The shadow of the Earth falls on the Moon during a lunar eclipse, and its computation involves calculating the \dev{च्छाया} (shadow), \dev{मध्यच्छाया} (middle shadow), and related quantities.

%% ========================================================================
\subsection{Verse 1: Beginning of the Chapter}

\vspace{0.5em}
\noindent\textbf{Sūtra (\dev{सूत्रम्}) — Verse 1:}

\begin{paracol}{2}
\backgroundcolor{c[0](.95\columnsep,.95\textheight)}{sanskritbg}
\backgroundcolor{c[1](.95\columnsep,.95\textheight)}{englishbg}

\begin{sanskritverse}
अथ चन्द्रच्छायाधिकारः प्रारभ्यते। तत्रादौ तदारम्भप्रयोजनमाह—
\end{sanskritverse}

\begin{sanskritverse}
प्राक् चन्द्रलग्नयोरन्तरं सप्तलग्नादीनां यतस्ततं पश्चात्।
प्रतिदिनमिन्दोश्च्छाया यतस्तदनयनमभिचार्ये॥ १ ॥
\end{sanskritverse}

\switchcolumn

\textbf{Introduction:} ``Now the chapter on the Moon's shadow is begun. First, the purpose of its commencement is stated ---''

\textbf{Translation:} ``Previously, the difference between the Moon and the \textit{lagna}, and afterwards, that of the \textit{saptalagna} etc., because from that the shadow of the Moon (\textit{indu-ccchāyā}) is daily [computed], therefore I shall state the computation (\textit{anayana}) of that.''

\end{paracol}

\vspace{0.5em}

\begin{paracol}{2}
\backgroundcolor{c[0](.95\columnsep,.95\textheight)}{sanskritbg}
\backgroundcolor{c[1](.95\columnsep,.95\textheight)}{englishbg}

\noindent\textbf{Sūtrārtha (\dev{सूत्रार्थः}):}
\begin{sanskritverse}
प्राक् (पूर्वस्थितिजे) चन्द्रलग्नयोरन्तरं, पश्चात् (पश्चिमस्थितिजे) सप्तलग्नचन्द्रयोः
अन्तरं च प्रतिदिनं यतनेन यतः (यस्मात्कारणात्) इन्दुच्छाया (चन्द्रच्छाया) ज्ञायते,
तथेन्दुच्छायातस्तदनन्तरं च ज्ञायते ततस्तदानयन (छायानयन) मभिचार्ये (कथयिष्ये) इति।
सिद्धान्तशेखरे ``प्रायतस्तु हीनरविमलग्नयोरन्तरलग्ननाशिनोश्च पश्चिमे'' इत्यादि
इनेन श्रीपतिनाचार्योक्तानुरूपमेवोक्तमिति ॥ १ ॥
\end{sanskritverse}

\switchcolumn

\textbf{Literal Meaning:} ``Previously (in the eastern horizon), the difference between the Moon and the \textit{lagna}, and afterwards (in the western horizon), the difference between the \textit{saptalagna} and the Moon, and daily by effort --- by which cause the Moon's shadow (\textit{indu-ccchāyā}) is known, and from the Moon's shadow that difference is [also] known --- therefore I shall state that computation (\textit{anayana}), the shadow-computation (\textit{chāyā-anayana}). In the \textit{Siddhāntaśekhara}: `Usually the difference between the diminished-Sun-\textit{lagna} and the \textit{lagna-nāśin} in the west' etc. --- by this, what was spoken by Śrīpati-nācārya is here spoken in the same manner.''

\end{paracol}

\vspace{0.5em}

\begin{paracol}{2}
\backgroundcolor{c[0](.95\columnsep,.95\textheight)}{sanskritbg}
\backgroundcolor{c[1](.95\columnsep,.95\textheight)}{englishbg}

\noindent\textbf{Hindi Bhāṣya (\dev{हिन्दी भाष्य}):}
\begin{sanskritverse}
अब चन्द्रच्छायाधिकार प्रारम्भ किया जाता है, पहले उसके प्रारम्भ-प्रयोजन को कहते हैं।
पूर्व स्थितिजे में चन्द्र और लग्न का अन्तर, पश्चिम स्थितिजे में सप्तलग्न और चन्द्र का अन्तर
प्रत्येक दिन में जो होता है उसी से क्योंकि चन्द्रच्छाया समझी जाती है (प्रथम चन्द्रच्छाया
का ज्ञान होता है) तथा चन्द्रच्छाया से उस अन्तर का ज्ञान होता है इसलिये छायानयन को
कहता हूँ॥ १ ॥
\end{sanskritverse}

\switchcolumn

\textbf{English Translation:} ``Now the chapter on the Moon's shadow is begun. First, its purpose is stated. The difference between the Moon and the \textit{lagna} in the eastern horizon, and the difference between the \textit{saptalagna} and the Moon in the western horizon --- because from these, daily, the Moon's shadow is known (first, the knowledge of the Moon's shadow is obtained), and from the Moon's shadow the knowledge of that difference is obtained --- therefore I state the shadow-computation.''

\end{paracol}

%% ========================================================================
\subsection{Verse 2: Shadow Computation}

\vspace{0.5em}
\noindent\textbf{Sūtra (\dev{सूत्रम्}) — Verse 2:}

\begin{paracol}{2}
\backgroundcolor{c[0](.95\columnsep,.95\textheight)}{sanskritbg}
\backgroundcolor{c[1](.95\columnsep,.95\textheight)}{englishbg}

\begin{sanskritverse}
प्रभहुण्डाकृतिकोर्ध्वकालिकोद्बलनखण्डसमस्थितिपाताः।
छन्दोव्यासस्तलग्ने स्वभग्रप्रायाण्न सङ्क्षेपः॥ २ ॥
\end{sanskritverse}

\switchcolumn

\textbf{Translation:} ``The parts of the upright-time-difference (\textit{ūrdhvakālika-udbala-naṣṭa-sama-sthiti-pāta}) [having] the form of a \textit{prabhahṛṇḍa} [spherical triangle], the semi-diameter (\textit{vyāsa}) of the shadow at the horizon (\textit{tala-lagna}), the proximity to its own \textit{agra} --- there is no abbreviation (\textit{saṅkṣepaḥ}).''

\textit{Note:} This is a technical verse whose exact interpretation requires the accompanying commentary. The \textit{prabhahṛṇḍa} is a type of spherical triangle used in Indian astronomy.

\end{paracol}

%% ========================================================================
\subsection{Verses 3--4: Detailed Shadow Calculation}

\vspace{0.5em}
\noindent\textbf{Sūtra (\dev{सूत्रम्}) — Verses 3--4:}

\begin{paracol}{2}
\backgroundcolor{c[0](.95\columnsep,.95\textheight)}{sanskritbg}
\backgroundcolor{c[1](.95\columnsep,.95\textheight)}{englishbg}

\begin{sanskritverse}
इदानीं कर्तव्यतामाह—
\end{sanskritverse}

\begin{sanskritverse}
प्राग्रहुण्डाकृतिकोर्ध्वकालिकोद्बलनखण्डसमस्थितिपाताः।
छन्दोव्यासस्तलग्ने स्वभग्रप्रायाण्न सङ्क्षेपः॥ ३ ॥
\end{sanskritverse}

\begin{sanskritverse}
मूलोनं मूलं गतं यद्वदुद्धृत्यागतं ततोऽन्यत्।
स्वभुक्त्या हृतमभीष्टच्छायागतिः सैव मूलोद्धृता॥ ४ ॥
\end{sanskritverse}

\switchcolumn

\textbf{Introduction to Verse 3:} ``Now [the author] states what is to be done ---''

\textbf{Translation — Verse 3:} [Same structure as verse 2, establishing the method.]

\textbf{Translation — Verse 4:} ``As the root (\textit{mūla}) diminished by the root, or the elapsed [time], having been extracted, another [quantity] comes therefrom. Divided by the own motion (\textit{svabhakti}), that is the desired shadow-motion (\textit{abhīṣṭa-ccchāyā-gati}). That same [quantity], extracted from the root (\textit{mūlodḍhṛtā}).''

\textit{Note:} This verse describes an iterative root-extraction method for computing the rate of change of the shadow. The \textit{mūla} here refers to a base quantity from which differences are taken.

\end{paracol}

\newpage


%% ========================================================================
%% CHAPTER 4: GRAHODAYĀSTĀDHIKĀRAḤ (Planetary Rising/Setting Chapter)
%% ========================================================================
\section{\texorpdfstring{{\devanagarifont ग्रहोदयास्ताधिकारः} — Grahodayāstādhikāraḥ (Planetary Rising and Setting)}{Grahodayāstādhikāraḥ (Planetary Rising and Setting)}}
\label{sec:grahodayasta}

\begin{center}
\textcolor{rulecolor}{\rule{0.8\textwidth}{0.5pt}}
\end{center}

This chapter (\dev{ग्रहोदयास्ताधिकारः}) deals with the computation of the rising and setting times of planets. The chapter explains how to determine the times when planets become visible (\dev{उदय}, heliacal rising) and invisible (\dev{अस्त}, heliacal setting) based on their elongation from the Sun.

%% ========================================================================
\subsection{Verse 1: Pratyaya (Introduction)}

\vspace{0.5em}
\noindent\textbf{Sūtra (\dev{सूत्रम्}) — Verse 1:}

\begin{paracol}{2}
\backgroundcolor{c[0](.95\columnsep,.95\textheight)}{sanskritbg}
\backgroundcolor{c[1](.95\columnsep,.95\textheight)}{englishbg}

\begin{sanskritverse}
अथोदयास्ताधिकारः प्रारभ्यते। तत्रादौ तदारम्भप्रयोजनमाह—
\end{sanskritverse}

\begin{sanskritverse}
ग्रहभास्करान्तरं प्राक् पद बाधकं ग्रहान्तरे यस्मात्।
स्वांशैर्यथाह्यास्तस्माद् द्युते तदानयनम्॥ १ ॥
\end{sanskritverse}

\switchcolumn

\textbf{Introduction:} ``Now the chapter on planetary rising and setting is begun. First, the purpose of its commencement is stated ---''

\textbf{Translation:} ``Because previously, in the inter-planetary computation, the \textit{grahabhāskara} distance is an obstruction (\textit{bādhaka}), therefore, with its own degrees, so that there may be agreement --- I state the computation of \textit{dyuti}.''

\textit{Note:} This verse serves as a parallel to the solar eclipse chapter's opening, establishing that Brahmagupta's method supersedes earlier approaches that failed to achieve agreement (\textit{aiṃya}) between computation and observation.

\end{paracol}

\vspace{0.5em}

\begin{paracol}{2}
\backgroundcolor{c[0](.95\columnsep,.95\textheight)}{sanskritbg}
\backgroundcolor{c[1](.95\columnsep,.95\textheight)}{englishbg}

\noindent\textbf{Sūtrārtha (\dev{सूत्रार्थः}):}
\begin{sanskritverse}
प्राक् पूर्वे क्षितिजे प्रभास्करान्तरैः स्वांशैः स्वकालांशैः पश्चात् पश्चिमक्षितिजे
ग्रहास्तरैः स्वकालांशैर्यथासंख्यं ग्रहाः हश्याहश्याः भवन्ति। तस्मात् तदानयनं
प्रहर्कान्तरानयनं वक्ष्य इति। अर्थात् प्राक्-क्षितिजे यस्मिन् दिने ग्रहोदयादनन्तरं
कालांशघटीभिः रवेरुदयस्तस्मिन् दिने प्रभो रात्रिशेषे हश्यो भवति। एवं यस्मिन् दिने
पश्चिमक्षितिजे रव्यस्तानन्तरं कालांशघटीमित्रं हस्यास्तमयस्तस्मिन् दिने सायंकाले
ग्रहस्य हश्यत्वम्॥ १ ॥
\end{sanskritverse}

\switchcolumn

\textbf{Literal Meaning:} ``Previously, in the eastern horizon, by the \textit{grahabhāskara} distances in their own degrees, in their own time-degrees, later in the western horizon, by the planetary distances in their own time-degrees, according to number, the planets become visible or invisible (\textit{haśya-āhaśya}). Therefore, I shall state that computation --- the computation of the \textit{praharāntara}. That is: in the eastern horizon, on whatever day after the rising of the planet, the rising of the Sun occurs after a certain number of time-ghaṭīs, on that day the planet becomes invisible (\textit{haśya}) in the remaining night. Similarly, in the western horizon, after the setting of the Sun, when the setting occurs after a certain number of time-ghaṭīs, on that day in the evening the planet's invisibility (\textit{haśyatva}) [is established].''

\end{paracol}

%% ========================================================================
\subsection{Verses 8--9: Daily Rising and Setting}

\vspace{0.5em}
\noindent\textbf{Sūtra (\dev{सूत्रम्}) — Verses 8--9:}

\begin{paracol}{2}
\backgroundcolor{c[0](.95\columnsep,.95\textheight)}{sanskritbg}
\backgroundcolor{c[1](.95\columnsep,.95\textheight)}{englishbg}

\begin{sanskritverse}
तात्कालिक ग्रह लग्न से असकृत् कर्म द्वारा प्रतिदिन उदय और अस्त साधन करना
प्रथम जिस समय ग्रहोदित घटी ज्ञात करना हो उस समय के लग्न (तात्कालिक लग्न)
तथा ग्रहोदय लग्न साधन कर दोनों के अन्तरघटी साधन कर उन (साधित घटी) से ग्रह
को चलाकर उसके फिर उदयलग्न साधन कर तात्कालिक स्थिर लग्न से पुनः इष्टघटी
साधन करना, इसतरह असकृत् (वारंवार) करना॥ ८ ॥
\end{sanskritverse}

\switchcolumn

\textbf{Translation — Verse 8:} ``By iterative computation (\textit{asakṛt-karma}) from the instantaneous (\textit{tātkālika}) planet-\textit{lagna}, the daily rising and setting are to be computed. First, at whatever time the \textit{grahodita-ghaṭī} [elapsed ghaṭīs since planetary rising] is to be known, the \textit{lagna} at that time (\textit{tātkālika lagna}) and the planetary rising-\textit{lagna} are to be computed. Computing the difference-ghaṭīs of both, by those (computed ghaṭīs) moving the planet, again computing its rising-\textit{lagna}, from the instantaneous fixed \textit{lagna} again computing the desired ghaṭīs --- in this manner, iteratively (\textit{asakṛt}, repeatedly) doing [this].''

\textit{Note:} This verse describes the iterative procedure (\textit{asakṛt-karma}) for refining the computation of a planet's rising time. The method involves successive approximation until the computed value stabilizes.

\end{paracol}

%% ========================================================================
\subsection{Verse 9: Moon's Rising and Setting Computation}

\vspace{0.5em}
\noindent\textbf{Sūtra (\dev{सूत्रम्}) — Verse 9:}

\begin{paracol}{2}
\backgroundcolor{c[0](.95\columnsep,.95\textheight)}{sanskritbg}
\backgroundcolor{c[1](.95\columnsep,.95\textheight)}{englishbg}

\begin{sanskritverse}
उदयास्तमयाविन्दोः कालांशैरंकसंमितैः कार्यः।
हीनत्वं त्वधिकत्वं तदन्तरे योगकालः स्यात्॥ ९ ॥
\end{sanskritverse}

\switchcolumn

\textbf{Translation:} ``The rising and setting of the Moon (\textit{indu}) are to be computed by time-degrees (\textit{kālāṃśa}) equal to digits (\textit{aṅka-saṃmitaiḥ}). If there is diminution (\textit{hīnatva}) or excess (\textit{adhikatva}), in that interval the \textit{yoga-kāla} [conjunction time] would be [obtained].''

\textit{Note:} The \textit{aṅka} (digit) is a unit equal to 12 degrees. The verse states that the Moon's rising and setting are computed using time-degrees measured in these units.

\end{paracol}

\vspace{0.5em}

\begin{paracol}{2}
\backgroundcolor{c[0](.95\columnsep,.95\textheight)}{sanskritbg}
\backgroundcolor{c[1](.95\columnsep,.95\textheight)}{englishbg}

\noindent\textbf{Sūtrārtha (\dev{सूत्रार्थः}):}
\begin{sanskritverse}
अंकसंमितैः कालांशैररात् द्वादशकालांशैरिन्दोश्चन्द्रस्योदयास्तमयौ साध्यौ।
अत्र पठितकालांशेन इष्टकालांशातां यदि हीनत्वं वाधिकत्वं भवेत् तदा तदन्तरे तयोः
पठितकालांशयोरन्तरे योगकालो योगवत्‌ वा योगवत्‌ कालः स्यादिति सप्तमश्लोकेन
स्फुटोक्तम्॥ ९ ॥
\end{sanskritverse}

\switchcolumn

\textbf{Literal Meaning:} ``By time-degrees equal to digits (\textit{aṅka}), that is, by twelve time-degrees (\textit{dvādaśa-kālāṃśa}), the rising and setting of the Moon (\textit{indu} = \textit{candra}) are to be computed. Here, if the computed time-degree (\textit{paṭhita-kālāṃśa}) has diminution or excess from the desired time-degree (\textit{iṣṭa-kālāṃśa}), then in that interval, between those two computed time-degrees, the \textit{yoga-kāla} (conjunction time), or \textit{yogavat-kāla} (time like conjunction), would be [obtained], as clearly stated in the seventh verse.''

\end{paracol}

\vspace{0.5em}

\begin{paracol}{2}
\backgroundcolor{c[0](.95\columnsep,.95\textheight)}{sanskritbg}
\backgroundcolor{c[1](.95\columnsep,.95\textheight)}{englishbg}

\noindent\textbf{Hindi Bhāṣya (\dev{हिन्दी भाष्य}):}
\begin{sanskritverse}
अंकसंमितैः कालांशैररथात् द्वादशतुल्यैः कालांशैरुदयास्तमयौ साध्यो। पठितकालांशस्य
इष्टकालांशस्य यदि हीनत्वं (अल्पत्वं) वा अधिकत्वं भवेत् तदन्तरे तयोः कालांशयोरन्तरे
योगकालः स्यात्॥ ९ ॥
\end{sanskritverse}

\switchcolumn

\textbf{English Translation:} ``By time-degrees equal to digits (\textit{aṅka}), that is, by time-degrees equal to twelve, the rising and setting are to be computed. If the computed time-degree has diminution (\textit{hīnatva}, being smaller) or excess (\textit{adhikatva}) from the desired time-degree, then in that interval, between those two time-degrees, the \textit{yoga-kāla} (conjunction time) would be [obtained].''

\end{paracol}

\newpage


%% ========================================================================
%% CHAPTER 5: ŚUKLONNATYADHIKĀRAḤ (Moon's Brightness Chapter)
%% ========================================================================
\section{\texorpdfstring{{\devanagarifont शुक्लोन्नत्यधिकारः} — Śuklonnatyadhikāraḥ (The Brightness of the Moon)}{Śuklonnatyadhikāraḥ (The Brightness of the Moon)}}
\label{sec:shuklonnati}

\begin{center}
\textcolor{rulecolor}{\rule{0.8\textwidth}{0.5pt}}
\end{center}

This chapter (\dev{शुक्लोन्नत्यधिकारः}) deals with the computation of the illuminated portion of the Moon's disk --- the \dev{शुक्ल} (bright) portion. The chapter explains how to determine the amount of the Moon's surface that is illuminated by the Sun at any given time.

%% ========================================================================
\subsection{Verses 1--3: Sphuṭaśuklonnati Computation}

\vspace{0.5em}
\noindent\textbf{Sūtra (\dev{सूत्रम्}) — Verse 4:}

\begin{paracol}{2}
\backgroundcolor{c[0](.95\columnsep,.95\textheight)}{sanskritbg}
\backgroundcolor{c[1](.95\columnsep,.95\textheight)}{englishbg}

\begin{sanskritverse}
अथ शुक्लोन्नत्यधिकारः प्रारभ्यते।
\end{sanskritverse}

\begin{sanskritverse}
रविचन्द्रपातलग्नैः स्वक्रान्त्युदयात् स्वलग्नगतशेषाः।
घटिकाः खचरार्धास्तात् स्वेष्टी रविशीतथू कृत्वा॥ ४ ॥
\end{sanskritverse}

\switchcolumn

\textbf{Introduction:} ``Now the chapter on the Moon's bright elevation (\textit{śuklonnati}) is begun.''

\textbf{Translation — Verse 4:} ``The ghaṭikās [which are] the \textit{khacarārdha} (half of the celestial sphere's motion) --- by the Sun, Moon, node, and \textit{lagna}, and from the rising of the own declination (\textit{sva-krānty-udaya}), [which are] the remaining [ghaṭikās] elapsed from the own \textit{lagna} --- having made [the computation] from the setting of the Sun, [these are] to be computed.''

\textit{Note:} The \textit{khacarārdha} refers to half the diurnal motion of the celestial sphere. The verse establishes the parameters needed to compute the Moon's illumination.

\end{paracol}

\vspace{0.5em}

\begin{paracol}{2}
\backgroundcolor{c[0](.95\columnsep,.95\textheight)}{sanskritbg}
\backgroundcolor{c[1](.95\columnsep,.95\textheight)}{englishbg}

\noindent\textbf{Sūtrārtha (\dev{सूत्रार्थः}):}
\begin{sanskritverse}
खचरार्धास्ताद् रव्यस्तात् स्वलग्नगतशेषाः शशिन उदयलग्नस्य गता दोषा वा
घटिकाः साध्याः। कैः। रविचन्द्रपातलग्नैस्तथा स्वक्रान्त्युदयात्। अत्रैतदुक्तं भवति।
रव्यस्तकाले रविचन्द्रः पातो लग्नम्। एतानि कृत्वा स्वक्रान्त्या चन्द्रक्रान्त्या
उदयात् चन्द्रोदयलग्नाज्ञोदयलग्नस्य गता वा दोषा घटिका ऊनस्य भोग्योद्धिकभुक्तियुक्तो
मध्योदयाद्य इत्यनेन विधिना साध्याः। सूर्यास्तानन्तरं यावतीभिः षटिकाभिश्चन्द्रास्तस्ताः
गता घटिकाः। सूर्यास्तात् प्राग् यावतीष्ठिकाभिश्चन्द्रास्तस्ता एषा घटिकाः। एवं प्राक्
क्षितिजे रव्युदयात् गता एषा वा चन्द्रोदयघटिकाः साध्या इत्यर्थादवगम्यते॥ ४ ॥
\end{sanskritverse}

\switchcolumn

\textbf{Literal Meaning:} ``The \textit{khacarārdha} ghaṭikās [which are] from the setting of the Sun (\textit{ravyastāt}), [which are] the remaining [ghaṭikās] elapsed from the own-\textit{lagna} of the Moon --- or the elapsed/deficient ghaṭikās of the Moon's rising-\textit{lagna} --- are to be computed. By what? By the Sun, Moon, node, \textit{lagna}, and from the rising of the own declination. Here, this is stated: at the time of the Sun's setting, the Sun, Moon, node, and \textit{lagna}. Having computed these, by the own declination (\textit{sva-krānti}), by the Moon's declination (\textit{candra-krānti}), from the rising, [computing] the elapsed or deficient ghaṭikās of the Moon's rising-\textit{lagna} from the \textit{udaya-lagna} by the method `diminished by the \textit{bhoga}, increased by the \textit{bhukti}, [gives] the middle \textit{udaya}', etc. After the Sun's setting, the ghaṭikās by which the Moon sets --- those are the elapsed ghaṭikās. Before the Sun's setting, the ghaṭikās by which the Moon sets --- those are the deficient ghaṭikās. Similarly, in the eastern horizon, from the Sun's rising, the elapsed or deficient ghaṭikās of the Moon's rising are to be computed --- this is to be understood from the meaning.''

\end{paracol}

\vspace{0.5em}

\begin{paracol}{2}
\backgroundcolor{c[0](.95\columnsep,.95\textheight)}{sanskritbg}
\backgroundcolor{c[1](.95\columnsep,.95\textheight)}{englishbg}

\noindent\textbf{Hindi Bhāṣya (\dev{हिन्दी भाष्य}):}
\begin{sanskritverse}
तात्कालिक ग्रह लग्न से असकृत् कर्म द्वारा प्रतिदिन उदय और अस्त साधन करना प्रथम
जिस समय ग्रहोदित घटी ज्ञात करना हो उस समय के लग्न (तात्कालिक लग्न) तथा
ग्रहोदय लग्न साधन कर दोनों के अन्तरघटी साधन कर उन (साधित घटी) से ग्रह को
चलाकर उसके फिर उदयलग्न साधन कर तात्कालिक स्थिर लग्न से पुनः इष्टघटी
साधन करना, इसतरह असकृत् (वारंवार) करना, एवं भस्तकाल ज्ञान के लिये भस्तलग्न (सप्तभस्त
लग्न) से कर्म करना चाहिये। इस तरह प्रथम चन्द्र दर्शन दिन से पूर्णिमा पर्यन्त इष्ट
दिन में चन्द्र के सूर्यास्त से त्रा सूर्योदय से अन्तमय काल और उदयिक काल साधन करना॥ ४ ॥
\end{sanskritverse}

\switchcolumn

\textbf{English Translation:} ``By iterative computation (\textit{asakṛt-karma}) from the instantaneous planet-\textit{lagna}, the daily rising and setting are to be computed. First, at whatever time the planetary elapsed ghaṭīs are to be known, the \textit{lagna} at that time (\textit{tātkālika lagna}) and the planetary rising-\textit{lagna} are to be computed. Computing the difference-ghaṭīs of both, by those (computed ghaṭīs) moving the planet, again computing its rising-\textit{lagna}, from the instantaneous fixed \textit{lagna} again computing the desired ghaṭīs --- in this manner, iteratively (\textit{asakṛt}, repeatedly) doing [this]. And for the knowledge of the \textit{sthita-kāla}, computation from the \textit{sthita-lagna} (\textit{sapta-bhastala-lagna}) should be done. In this way, from the first day of lunar visibility (\textit{candra darśana}) up to the full moon (\textit{pūrṇimā}), on the desired day, the setting time from the Sun's setting and the rising time from the Sun's rising are to be computed.''

\end{paracol}

%% ========================================================================
\subsection{Verse 5: The {\devanagarifont शुक्लोन्नति} (Bright Elevation)}

\vspace{0.5em}
\noindent\textbf{Sūtra ({\devanagarifont सूत्रम्}) — Verse 5:}

\begin{paracol}{2}
\backgroundcolor{c[0](.95\columnsep,.95\textheight)}{sanskritbg}
\backgroundcolor{c[1](.95\columnsep,.95\textheight)}{englishbg}

\begin{sanskritverse}
रविचन्द्रपातलग्नैः स्वक्रान्त्युदयात् स्वलग्नगतशेषाः।
घटिकाः खचरार्धास्तात् स्वेष्टी रविशीतथू कृत्वा॥ ५ ॥
\end{sanskritverse}

\switchcolumn

\textbf{Translation — Verse 5:} [Recapitulates the parameters for computing the bright elevation.]

\textit{Technical summary:} The verse reiterates that by the Sun, Moon, node, \textit{lagna}, and from the rising of the declination, the \textit{khacarārdha} ghaṭikās from the Sun's setting are to be used to compute the Moon's bright elevation.

\end{paracol}

\vspace{0.5em}

\begin{paracol}{2}
\backgroundcolor{c[0](.95\columnsep,.95\textheight)}{sanskritbg}
\backgroundcolor{c[1](.95\columnsep,.95\textheight)}{englishbg}

\noindent\textbf{Upapatti (\dev{उपपत्तिः}):}
\begin{sanskritverse}
यस्मिन् दिने सूर्यास्तादनन्तरं कालांशघटीभोद्धिकाभिश्चन्द्रास्तः। रव्युदयात् प्राक्
कालांशघटीभोद्धिकाभिर्चन्द्रोदयस्तस्मिन्नेव चन्द्रदर्शनं सति द्वे श्वज्योन्नतिः साध्या।
सितशुक्लोन्नतिरच बिम्बार्धं चन्द्रसितम् एवं भास्करेण मासान्तपादे प्रथमे च
शुक्लोन्नतिः साधिता। द्वितीयतृतीयपादयोर्बिम्बार्धं कृष्णं भवति। ततस्तत्र
कृष्णशुक्लोन्नतिर्भवति। इत्याचार्येण शुक्लोन्नतिद्वयं साध्यते। तदयं पादचर्चा न
कृता। भास्करेण न हि स्पष्टा कृष्णशुक्लोन्नतिः प्रेक्ष्यते इति मनसि संप्रधार्य
सितशुक्लोन्नतिरेव साधिता। उक्तं च ``द्वितीयतृतीययोरपि चरणयोर्ब्रह्मगुप्तादयः कृष्णशुक्लोन्नतिरानीता सा मम न संमता। न हि नरे कृष्णः शुक्लोन्नतिः स्पष्टोपलक्ष्यते''॥
\end{sanskritverse}

\switchcolumn

\textbf{English Translation of Proof:} ``On whatever day, after the Sun's setting, the Moon sets after an additional number of time-ghaṭīs. Before the Sun's rising, the Moon rises [after a certain number of ghaṭīs]. On that very day, when the Moon is visible, the two bright elevations (\textit{śva-jy-unnati}) are to be computed. The bright elevation of the bright [half] (\textit{sita-śuklonnati}) is half the disk (\textit{bimbārdha}) of the Moon's bright [portion]. Thus Bhāskara computed the \textit{śuklonnati} in the first \textit{pāda} of the month-end. In the second and third \textit{pāda}s, half the disk is dark (\textit{kṛṣṇa}). Therefore there the dark \textit{śuklonnati} occurs. Thus by the teacher (\textit{ācārya}), two \textit{śuklonnati}s are computed. But here, the \textit{pāda} discussion is not made. Bhāskara, having considered in his mind that the dark \textit{śuklonnati} is not clearly observable, computed only the bright \textit{śuklonnati}. And it is stated: `Brahmagupta and others have produced the dark \textit{śuklonnati} also for the second and third \textit{caraṇa}s; that is not approved by me. For neither is the dark \textit{śuklonnati} clearly perceived by men.'

\end{paracol}

%% ========================================================================
\subsection{Verses 11--12: {\devanagarifont स्पष्टशुक्ल} (True Bright Portion) Computation}

\vspace{0.5em}
\noindent\textbf{Sūtra ({\devanagarifont सूत्रम्}) — Verses 11--12:}

\begin{paracol}{2}
\backgroundcolor{c[0](.95\columnsep,.95\textheight)}{sanskritbg}
\backgroundcolor{c[1](.95\columnsep,.95\textheight)}{englishbg}

\begin{sanskritverse}
तच्छाया गुणिते वा मध्यगती तिथिगुणस्वकरं हते।
फलयोदिक् साम्येऽन्तरमवनतिरंक्यं दिगन्यत्वे॥ ११ ॥
\end{sanskritverse}

\switchcolumn

\textbf{Translation — Verse 11:} ``Or when the middle motion (\textit{madhyagati}) is multiplied by its own shadow (\textit{tacchāyā}) and by the semi-diameter (\textit{sva-kara}) multiplied by the \textit{tithi} --- the results are additive or subtractive. In equality (\textit{sāmya}), the difference is [taken]; in direction-difference (\textit{dig-anyatva}), [it is] the \textit{avanati} [depression] expressed in digits (\textit{aṅkya}).''

\textit{Note:} This verse gives the method for computing the true bright portion (\textit{spaṣṭa-śukla}) of the Moon. The \textit{madhyagati} is the mean daily motion, and the \textit{tacchāyā} is the shadow corresponding to the planet's altitude.

\end{paracol}

\vspace{0.5em}

\begin{paracol}{2}
\backgroundcolor{c[0](.95\columnsep,.95\textheight)}{sanskritbg}
\backgroundcolor{c[1](.95\columnsep,.95\textheight)}{englishbg}

\noindent\textbf{Sūtrārtha (\dev{सूत्रार्थः}):}
\begin{sanskritverse}
हृज्ये रविचन्द्रयोः क्षेपी स्वस्वमध्यगतिगुरो तिथिगुणितव्यासदलेन पञ्चदशगुरिणा
त्रिज्यया भक्त हृते पृथक् पृथक् रवि-चन्द्र-योर्-नती भवतः। वा तयोर्मध्यगती
तयोः क्षेपसमं हृज्ययोः छाये ताम्यां गुरिते तिथिगुणस्वस्वच्छायाकर्महते तयोरन्तरी भवतः।
\end{sanskritverse}

\switchcolumn

\textbf{Literal Meaning:} ``In the \textit{hṛjya} [radius = 3438'] of the Sun and Moon, the \textit{kṣepī} [sine of latitude], having been multiplied by the respective middle motion and by the semi-diameter (\textit{vyāsadala}) multiplied by the \textit{tithi} and by fifteen [i.e., by the \textit{tithi}-multiplied semi-diameter multiplied by 15], divided by the \textit{trijyā} [radius = 3438'], separately, separately, the \textit{natī} [depressions] of the Sun and Moon are obtained. Or, their middle motion, their \textit{kṣepa}-equal \textit{hṛjya}, the shadow being multiplied by that, multiplied by the \textit{tithi}-multiplied own shadow-\textit{karman}, their difference is [taken].''

\end{paracol}

\vspace{0.5em}

\begin{paracol}{2}
\backgroundcolor{c[0](.95\columnsep,.95\textheight)}{sanskritbg}
\backgroundcolor{c[1](.95\columnsep,.95\textheight)}{englishbg}

\noindent\textbf{Upapatti (\dev{उपपत्तिः}):}
\begin{sanskritverse}
प्रथमानयनस्य ``सहक्षेपाघ्ना इन्दोर्निजमध्यभुक्तितिथ्यंशनिघ्ना'' इत्यादि
भास्करविधिना स्फुटा।
\end{sanskritverse}

\switchcolumn

\textbf{English Translation of Proof:} ``The first computation [is made] accurate by the method of Bhāskara: `Having multiplied by the \textit{sahakṣepa} [combined latitude] of the Moon, by the own middle motion-\textit{tithi}-aṃśa-multiplied [quantity]' etc.''

\end{paracol}

\newpage


%% ========================================================================
%% CHAPTER 6: CHANDRAŚUKLONNATYADHIKĀRAḤ (Moon's Bright Elevation, Chapter VII)
%% ========================================================================
\section{\texorpdfstring{{\devanagarifont चन्द्रशुक्लोन्नत्यधिकारः} — Chandraśuklonnatyadhikāraḥ (Moon's Bright Elevation)}{Chandraśuklonnatyadhikāraḥ (Moon's Bright Elevation)}}
\label{sec:chandra-shuklonnati}

\begin{center}
\textcolor{rulecolor}{\rule{0.8\textwidth}{0.5pt}}
\end{center}

This is the seventh \dev{अध्याय} (chapter) of the \dev{ब्राह्मस्फुटसिद्धान्त}, dealing with the Moon's bright elevation (\dev{चन्द्रशुक्लोन्नति}). It comprises eighteen verses (\dev{श्लोक}) and covers the computation of the arc of illumination on the Moon's disk.

%% ========================================================================
\subsection{Verses 1--2: Introduction}

\vspace{0.5em}
\noindent\textbf{Sūtra (\dev{सूत्रम्}) — Verses 1--2:}

\begin{paracol}{2}
\backgroundcolor{c[0](.95\columnsep,.95\textheight)}{sanskritbg}
\backgroundcolor{c[1](.95\columnsep,.95\textheight)}{englishbg}

\begin{sanskritverse}
अथ चन्द्रशुक्लोन्नत्यधिकारः प्रारम्भः।
\end{sanskritverse}

\begin{sanskritverse}
प्राग्रहुण्डाकृतिकोर्ध्वकालिकोद्बलनखण्डसमस्थितिपाताः।
छन्दोव्यासस्तलग्ने स्वभग्रप्रायाण्न सङ्क्षेपः॥ १ ॥
\end{sanskritverse}

\begin{sanskritverse}
मूलोनं मूलं गतं यद्वदुद्धृत्यागतं ततोऽन्यत्।
स्वभुक्त्या हृतमभीष्टच्छायागतिः सैव मूलोद्धृता॥ २ ॥
\end{sanskritverse}

\switchcolumn

\textbf{Introduction:} ``Now the beginning of the chapter on the Moon's bright elevation (\textit{candra-śukla-unnati}).''

\textbf{Translation — Verse 1:} ``The parts of the upright-time-difference (\textit{ūrdhvakālika-udbala-naṣṭa-sama-sthiti-pāta}) [having] the form of a \textit{prāgrahuṇḍa} [spherical triangle], the semi-diameter (\textit{vyāsa}) of the shadow at the horizon (\textit{tala-lagna}), the proximity to its own \textit{agra} --- there is no abbreviation (\textit{saṅkṣepaḥ}).''

\textbf{Translation — Verse 2:} ``As the root (\textit{mūla}) diminished by the root, or the elapsed [time], having been extracted, another [quantity] comes therefrom. Divided by the own motion (\textit{svabhakti}), that is the desired shadow-motion (\textit{abhīṣṭa-ccchāyā-gati}). That same [quantity], extracted from the root (\textit{mūlodḍhṛtā}).''

\textit{Note:} Verse 2 establishes an iterative root-extraction algorithm for computing the rate of change of the shadow, which is essential for determining the Moon's illumination over time.

\end{paracol}

%% ========================================================================
\subsection{Verse 18: Chapter Conclusion}

\vspace{0.5em}
\noindent\textbf{Sūtra (\dev{सूत्रम्}) — Verse 18:}

\begin{paracol}{2}
\backgroundcolor{c[0](.95\columnsep,.95\textheight)}{sanskritbg}
\backgroundcolor{c[1](.95\columnsep,.95\textheight)}{englishbg}

\begin{sanskritverse}
इदानीमध्यायोपसंहारमाह—
\end{sanskritverse}

\begin{sanskritverse}
इति बाहुकोटिकर्णस्फुटसितपरिलेखसूत्रकर्दमेषु।
द्वाराष्टादशा चन्द्रशुक्लोन्नतिसप्तमोऽध्यायः॥ १८ ॥
\end{sanskritverse}

\switchcolumn

\textbf{Introduction:} ``Now [the author] states the conclusion of the chapter ---''

\textbf{Translation — Verse 18:} ``Thus, in [the topics of] \textit{bāhu} (base), \textit{koṭi} (perpendicular), \textit{karṇa} (hypotenuse), \textit{spaṣṭa} (true [longitude]), \textit{sita} (bright [part]), \textit{parilekha} (outline), \textit{sūtra} (line), and \textit{kardama} [mud] --- by eighteen doors (\textit{dvāra}), [this is] the seventh chapter [on] the Moon's bright elevation (\textit{candra-śuklonnati}).''

\textit{Note:} The chapter lists eight topics (\textit{bāhu}, \textit{koṭi}, \textit{karṇa}, \textit{spaṣṭa}, \textit{sita}, \textit{parilekha}, \textit{sūtra}, \textit{kardama}) covered in 18 verses. The \textit{kardama} refers to the ``mud'' or penumbral shadow region.

\end{paracol}

\vspace{0.5em}

\begin{paracol}{2}
\backgroundcolor{c[0](.95\columnsep,.95\textheight)}{sanskritbg}
\backgroundcolor{c[1](.95\columnsep,.95\textheight)}{englishbg}

\noindent\textbf{Sūtrārtha (\dev{सूत्रार्थः}):}
\begin{sanskritverse}
परिलेखसूत्रमेव कर्ण इति... परिलेखसूत्रकर्णः। कोषं स्पष्टार्थम्॥ १८ ॥
\end{sanskritverse}

\switchcolumn

\textbf{Literal Meaning:} ``The \textit{parilekha-sūtra} itself [is] the \textit{karṇa}... the \textit{parilekha-sūtra-karṇaḥ}. The meaning is clear.''

\end{paracol}

\vspace{0.5em}

\begin{paracol}{2}
\backgroundcolor{c[0](.95\columnsep,.95\textheight)}{sanskritbg}
\backgroundcolor{c[1](.95\columnsep,.95\textheight)}{englishbg}

\noindent\textbf{Hindi Bhāṣya (\dev{हिन्दी भाष्य}):}
\begin{sanskritverse}
इस तरह बाहु, कोटि, कर्ण, स्फुट, सिताङ्गुल, परिलेखसूत्र (कर्ण) में अष्टादश
भावार्थश्लोकों से चन्द्रशुक्लोन्नति नामक सप्तम अध्याय है॥ १८ ॥
\end{sanskritverse}

\switchcolumn

\textbf{English Translation:} ``In this way, in [the topics of] \textit{bāhu}, \textit{koṭi}, \textit{karṇa}, \textit{spaṣṭa}, \textit{sitāṅgula}, \textit{parilekha-sūtra} (\textit{karṇa}), by eighteen meaning-verses, is the seventh chapter named \textit{Candraśuklonnati}.''

\textit{Note:} The \textit{sitāṅgula} (bright digit) is a unit for measuring the illuminated portion of the Moon's disk. One \textit{aṅgula} (digit) = 1/12 of the Moon's diameter.

\end{paracol}

\vspace{1em}
\begin{center}
\textcolor{rulecolor}{\rule{0.8\textwidth}{0.5pt}}
\end{center}

\subsection*{Concluding Verse ({\devanagarifont अन्तिमश्लोकः})}

\begin{paracol}{2}
\backgroundcolor{c[0](.95\columnsep,.95\textheight)}{sanskritbg}
\backgroundcolor{c[1](.95\columnsep,.95\textheight)}{englishbg}

\begin{sanskritverse}
मधुसूदनसूनुनोदितो यस्तिलकः श्रीपथुनेह जिष्णुजोक्तः।
हृदि तं विनिधाय नूतनोऽयं रचितः सूर्यगतिः सुधाकरेण॥
\end{sanskritverse}

\begin{sanskritverse}
इति श्रीकृपालुदत्तसूनुसुधाकरद्विवेदिविरचिते ब्राह्मस्फुटसिद्धान्तनुतनतिलके
चन्द्रशुक्लोन्नत्यधिकारः सप्तमः॥ ७० ॥
\end{sanskritverse}

\switchcolumn

\textbf{Translation — Concluding Verse:} ``That \textit{tilaka} (ornament) which was spoken by the son (\textit{sūnu}) of Madhusūdana, which was spoken here by Śrīpathi for the conqueror (\textit{jiṣṇu}) --- having placed that in [my] heart, this new [work] \textit{Sūryagati} has been composed by Sudhākara.''

\textbf{Colophon:} ``Thus, in the \textit{Brāhmasphuṭasiddhānta-nutana-tilaka} composed by Śrī Kṛpāludatta-sūnu Sudhākara Dvivedī, the seventh chapter --- the \textit{Candraśuklonnatyadhikāraḥ} --- [is concluded].''

\end{paracol}

\begin{center}
\textcolor{rulecolor}{\rule{0.8\textwidth}{0.5pt}}
\end{center}

\noindent\textbf{\dev{॥ इति श्रीब्रह्मगुप्तविरचिते ब्राह्मस्फुटसिद्धान्ते चन्द्रशुक्लोन्नत्यधिकारः सप्तमः समाप्तः ॥}}

\noindent\textbf{[Thus ends the Seventh Chapter --- the Moon's Bright Elevation --- in the \textit{Brāhmasphuṭasiddhānta} composed by Śrī Brahmagupta.]}

\vspace{2cm}

%% ========================================================================
%% APPENDIX: NOTES ON THE TEXT AND COMMENTARY
%% ========================================================================
\newpage
\section*{Appendix: Notes on the Text and Commentary}
\addcontentsline{toc}{section}{Appendix: Notes on the Text and Commentary}

\subsection*{About the {\devanagarifont ब्राह्मस्फुटसिद्धान्त}}

The {\devanagarifont ब्राह्मस्फुटसिद्धान्त} (\textit{Brāhmasphuṭasiddhānta}, ``Correctly Established Doctrine of Brahma'') is the magnum opus of the great Indian mathematician and astronomer \textbf{Brahmagupta} (598--668 CE). Composed in 628 CE at Bhinmal (in present-day Rajasthan, India), it consists of 24 chapters ({\devanagarifont अध्याय}) containing a total of 1008 verses ({\devanagarifont श्लोक}) in {\devanagarifont आर्या} and {\devanagarifont अनुष्टुभ्} metres.

\subsection*{Contents of This Chunk (Pages 281--420)}

This typeset volume covers the following chapters from Part II of the text:

\begin{enumerate}
\item \textbf{\dev{चन्द्रग्रहणविकारः}} --- Lunar Eclipse Variations \\
  Deals with the computation of eclipse magnitudes, durations, and the \textit{iṣṭagrāsa} (desired obscuration) method.

\item \textbf{\dev{सूर्यग्रहणाधिकारः}} --- Solar Eclipse Chapter \\
  Covers the computation of solar eclipses, including the five \textit{jyā} (sines) required: \textit{udayajyā}, \textit{madhyajyā}, \textit{ravidayākarṇa}, \textit{viśimaśaṅku}, and \textit{hṛkṣepa}.

\item \textbf{\dev{चन्द्रच्छायाधिकारः}} --- Moon's Shadow Chapter \\
  Treats the computation of the Earth's shadow as seen on the Moon's disk.

\item \textbf{\dev{ग्रहोदयास्ताधिकारः}} --- Planetary Rising and Setting \\
  Explains the methods for computing the heliacal rising (\textit{udaya}) and setting (\textit{asta}) of planets.

\item \textbf{\dev{शुक्लोन्नत्यधिकारः}} --- The Brightness of the Moon \\
  Computes the illuminated portion of the Moon's disk.

\item \textbf{\dev{चन्द्रशुक्लोन्नत्यधिकारः}} --- Moon's Bright Elevation \\
  The seventh chapter, with 18 verses, covering the arc of illumination.
\end{enumerate}

\subsection*{The Commentary ({\devanagarifont नूतनतिलक})}

The Hindi commentary ({\devanagarifont टीका}) in this edition was composed by \textbf{Paṇḍita Kṛpāludatta Sūnu Sudhākara Dvivedī}. The commentary follows the traditional structure of Indian mathematical commentaries:

\begin{itemize}
\item {\devanagarifont सूत्र} (\textit{Sūtra}) --- The original Sanskrit verse
\item {\devanagarifont सूत्रार्थ} (\textit{Sūtrārtha}) --- Literal meaning of the verse (marked {\devanagarifont सु. भा.})
\item \dev{विवृति} (\textit{Vivṛti}) --- Detailed explanation (marked \dev{वि. भा.})
\item \dev{हिन्दी भाष्य} (\textit{Hindi Bhāṣya}) --- Hindi explanation (marked \dev{हि. भा.})
\item \dev{उपपत्ति} (\textit{Upapatti}) --- Mathematical proof or demonstration
\end{itemize}

\subsection*{Mathematical Notation}

Brahmagupta's work represents one of the high points of Indian mathematics. Key contributions in these chapters include:

\begin{itemize}
\item Rules for computing eclipse magnitudes using the \dev{भुज} (\textit{bhuja}, sine) and \dev{कोटि} (\textit{koṭi}, cosine)
\item The \dev{कर्ण} (\textit{karṇa}) formula: $\text{karṇa} = \sqrt{\text{bhuja}^2 + \text{koṭi}^2}$
\item Iterative methods (\dev{असकृत्}, \textit{asakṛt}) for refining planetary positions
\item Proportional reasoning (\dev{त्रैराशिक}, \textit{trairāśika}) for time conversions
\item Geometric constructions using \dev{क्षेत्र} (\textit{kṣetra}, diagrams)
\end{itemize}

\subsection*{IAST Transliteration Conventions}

This bilingual edition uses the International Alphabet of Sanskrit Transliteration (IAST) scheme:

\begin{center}
\begin{tabular}{llllll}
\toprule
\textbf{Devanagari} & \textbf{IAST} & \textbf{Devanagari} & \textbf{IAST} & \textbf{Devanagari} & \textbf{IAST} \\
\midrule
अ & a & क & ka & त & ta \\
आ & ā & ख & kha & थ & tha \\
इ & i & ग & ga & द & da \\
ई & ī & घ & gha & ध & dha \\
उ & u & ङ & ṅ & न & na \\
ऊ & ū & च & ca & प & pa \\
ए & e & छ & cha & फ & pha \\
ऐ & ai & ज & ja & ब & ba \\
ओ & o & झ & jha & भ & bha \\
औ & au & ञ & ñ & म & ma \\
ऋ & ṛ & ट & ṭa & य & ya \\
ॄ & ṝ & ठ & ṭha & र & ra \\
\bottomrule
\end{tabular}
\end{center}

\vspace{1cm}
\begin{center